\begin{textsample}{NEG Dim 1 – 2012\_c6 – Score -1.00 – 2012\_c6/enade\_2012\_publi\_bc28e23025.txt}  \label{ex:f1_neg_260}
CAPÍTULO 3 ANÁLISE TÉCNICA DA PROVA Este capítulo tem por objetivo apresentar o desempenho dos estudantes concluintes de Comunicação Social - Publicidade e Propaganda no ENADE/2012.
Para isso, foram calculadas as estatísticas básicas da prova em seu todo, bem como as estatísticas dos componentes relacionadas à Formação Geral, ao de Conhecimento Específico da Área e das questões discursivas isoladamente.
Nas tabelas, são apresentados o tamanho da população inscrita e de presentes, e as seguintes estatísticas das notas9 : média do desempenho na prova, erro padrão da média, desvio padrão, nota mínima, mediana e nota máxima.
Tais estatísticas contemplam o total de estudantes concluintes da Área de Comunicação Social - Publicidade e Propaganda presentes à prova do ENADE 2012, tendo -se em vista as seguintes agregações : ( a ) as Grandes Regiões e o país como um todo ; ( b ) a Categoria Administrativa ; e ( c ) a Organização Acadêmica.
Em relação aos gráficos de distribuição de notas, o intervalo considerado foi de 10 unidades, aberto à esquerda e fechado à direita, com exceção do primeiro intervalo, [ 0 ; 10 ], fechado em ambos os extremos.
Para os gráficos de distribuição das notas das questões discursivas, foram consideradas mais duas categorias : questão em branco e nota zero. 3.1 ESTATÍSTICAS BÁSICAS DA PROVA 3.1.1 Estatísticas Básicas Gerais A Tabela 3.1 apresenta as estatísticas básicas da prova por grande Região.
A população total de inscritos foi de 19.733.
Destes, 15.239 \textbf{estiveram} presentes, sendo 22,8 \% o índice de não comparecimento.
A região de maior abstenção foi a Norte ( 37,7 \% ) e a de menor abstenção foi a Sul ( 14,4 \% ).
A média das notas da prova como um todo ( nas seções seguintes serão analisados os componentes de Formação Geral e de Conhecimento Específico ) foi 37,8, sendo que os alunos da região Centro-Oeste obtiveram a média mais baixa ( 36,0 ) e os da região Sul obtiveram a média mais alta ( 41,2 ).
As demais médias foram : 39,3 na região Nordeste e 37,0 nas regiões Norte e Sudeste.
O desvio padrão para o Brasil como um todo foi 13,8, sendo o maior desvio padrão encontrado nas regiões Nordeste, Sul e Centro-Oeste ( 13,9 ) e o menor na região Norte ( 13,0 ), indicando uma menor dispersão das notas desta última região.
A região que obteve a maior nota máxima foi a Nordeste ( 81,9 ), ao passo que a região que atingiu a menor nota máxima foi a Norte ( 76,8 ).
A mediana do Brasil como um todo foi 38,0, sendo a maior mediana obtida na região Sul ( 42,1 ) e a menor obtida na Centro-Oeste ( 35,9 ).
A nota mínima foi zero em todas as regiões, sem exceção.
O comportamento das notas dos estudantes de todo o Brasil pode ser observado no Gráfico 3.1 que apresenta um histograma com a distribuição das mesmas.
Essa é uma distribuição unimodal com moda no intervalo ( 30 ; 40 ].
Apesar do coeficiente de assimetria da distribuição das notas ser negativo ( -0,08 ) este é pequeno e podemos considerar que a distribuição é aproximadamente simétrica.
As distribuições por Grande Região também apresentam assimetria negativa, concentração pouco maior do lado direito do histograma e mais espalhada do lado esquerdo.
A única exceção é a região Centro-Oeste, com coeficiente de assimetria positivo ( 0,03 ), com o comportamento oposto : concentração à esquerda e cauda um pouco mais longa à direita.
Os Gráficos 3.2, 3.3 e 3.4 apresentam informações referentes à média da nota final dos Participantes, desagregando os resultados de acordo com, respectivamente, as Grandes Regiões do país, a Categoria Administrativa e a Organização Acadêmica.
Os gráficos apresentam o valor da média das notas como uma barra e os extremos do intervalo de confiança de 95 \% como linhas verticais unidas por uma linha horizontal na forma da letra H maiúscula.
Considerando -se o gráfico de notas segundo Grande Região ( Gráfico 3.2 ), observa -se que existe diferença estatisticamente significativa ao nível de 95 \% entre a maior média, obtida na região Sul ( 41,2 ) e a menor, obtida na região Centro-Oeste ( 36,0 ).
Levando -se em conta os estudantes por Categorias Administrativas ( Gráfico 3.3 ), observa -se que existe diferença estatisticamente significante entre as médias das notas das IES Públicas e Privadas.
Pode ser observado ainda que a média dos alunos de IES Públicas ( 41,0 ) é maior que a dos alunos de IES Privadas ( 37,6 ).
A diferença entre IES Públicas e Privadas ( 3,4 ) é menor do que a diferença entre regiões Sul e Centro-Oeste ( 5,2 ), a maior e a menor média entre as regiões.
Tendo como foco o Gráfico 3.4, que apresenta as notas médias das provas segundo Organização Acadêmica, constata -se que existe diferença estatisticamente significativa ao nível de 95 \% nas médias das notas dos estudantes provenientes de Universidades com os de Centros Universitários e Faculdades.
A maior média foi obtida pelos estudantes de Universidades ( 38,6 ), e a menor, pelos de Centros Universitários ( 36,7 ). 3.1.2 Estatísticas Básicas no Componente de Formação Geral A Tabela 3.2 apresenta as estatísticas básicas em relação ao componente da prova que avalia a Formação Geral dos estudantes concluintes.
Os alunos de todo Brasil obtiveram desempenho médio de 40,7.
Quanto à variabilidade, o desvio padrão das notas dos estudantes do Brasil como um todo foi 17,1.
A maior média foi obtida na região Sul ( 42,7 ), e a menor, na região Centro-Oeste ( 38,5 ).
As demais médias foram : 41,4 na região Nordeste, 40,7 na região Norte e 40,4 na região Sudeste.
Já o maior desvio padrão foi obtido na região Centro-Oeste ( 17,3 ) e o menor na região Sul ( 16,6 ).
Os demais desvios padrões foram : 17,2 na região Nordeste e 17,1 nas regiões Norte e Sudeste.
A maior nota no componente de Formação Geral da prova do ENADE foi obtida por pelo menos um aluno da região Sudeste ( 93,0 ) enquanto que a menor nota máxima foi obtida na região Centro-Oeste ( 83,5 ).
Nas outras regiões as notas máximas foram : 89,0 na região Sul, 88,0 na região Nordeste e 85,5 na região Norte.
A mediana do Brasil como um todo foi 41,5, sendo a menor mediana encontrada na região Centro-Oeste ( 38,5 ) e a maior encontrada na região Sul ( 43,5 ).
A nota mínima nesta parte foi zero em todas as regiões, sem exceção.
O Gráfico 3.5 propicia a avaliação do desempenho dos estudantes no componente de Formação Geral a partir do histograma da distribuição das notas correspondentes.
A distribuição é unimodal, com moda em ( 40 ; 50 ], enquanto na prova como um todo a moda foi alcançada no intervalo ( 30 ; 40 ].
Nota -se, ainda, que no gráfico 3.5 as notas apresentam uma maior dispersão do que no Gráfico 3.1 ( distribuição das notas da prova ), confirmado pela comparação dos desvios padrões : 13,8 para a nota da prova como um todo e 17,1 para o componente de Formação Geral.
Para o componente de Formação Geral, o coeficiente de assimetria da distribuição das notas dos estudantes, como na prova como um todo, também é negativo ( -0,20 ).
Por ser um coeficiente pequeno, a distribuição apresenta uma leve concentração à direita e cauda maior à esquerda.
Em todas as Grandes Regiões os histogramas também possuem assimetria negativa ( entre -0,31 para a região Sul e -0,10 para a região Norte ).
Nos Gráficos 3.6, 3.7 e 3.8 são apresentadas as informações referentes ao desempenho dos Concluintes no componente de Formação Geral, em diferentes agregações : Grande Região do país, Categoria Administrativa e Organização Acadêmica.
Observa -se pelo Gráfico 3.6 que existe diferença estatisticamente significativa entre a maior e a menor média das notas no Componente de Formação Geral, obtidas nas regiões Sul ( 42,7 ) e Centro-Oeste ( 38,5 ) respectivamente.
Vemos que o intervalo de confiança mais largo é o da região Norte ; já o intervalo mais estreito é observado na região Sudeste.
Este fato \textbf{está} relacionado, também, com o tamanho da população envolvida, menor na região Norte do que na Sudeste.
No Gráfico 3.7, que representa as notas médias no Componente de Formação Geral segundo Categoria Administrativa do país, observa -se que não existe diferença estatisticamente significativa entre as médias.
Ao contrário do que ocorreu para as notas da prova como um todo, para o componente de Formação Geral, os concluintes de Comunicação Social - Publicidade e Propaganda, das IES Públicas ( 40,6 ) obtiveram uma média menor do que os das IES Privadas ( 40,7 ).
Considerando -se o tipo de Organização Acadêmica, nota -se, no Gráfico 3.8, uma diferença estatisticamente significativa entre a maior e a menor média.
Nas Universidades ( 41,2 ), a média é mais elevada do que em Faculdades ( 40,3 ) e Centros Universitários ( 39,8 ). 3.1.3 Estatísticas Básicas do Componente de Conhecimento Específico A Tabela 3.3 apresenta as estatísticas básicas referentes ao componente de Conhecimento Específico da área de Comunicação Social - Publicidade e Propaganda.
A média do desempenho dos alunos do Brasil como um todo foi 36,9.
A maior média foi obtida na região Sul ( 40,7 ), e a menor, na região Centro-Oeste ( 35,2 ).
As demais médias foram : 38,6 na região Nordeste, 35,9 na região Sudeste e 35,7 na região Norte.
Quanto à variabilidade das notas, o desvio padrão do Brasil como um todo foi 14,9, sendo o maior desvio padrão observado na região Sul ( 15,2 ) e o menor na região Norte ( 13,9 ).
Os demais desvios foram : 14,9 da região Nordeste, 14,8 da região Centro-Oeste e 14,7 da região Sudeste.
A mediana das notas dos estudantes de todo o Brasil foi 36,7.
A maior mediana ocorreu na região Sul ( 41,4 ) e a menor nas regiões Sudeste e Centro-Oeste ( 35,4 ).
As demais medianas foram 38,9 na região Nordeste e 35,7 na região Norte.
A nota máxima do Brasil como um todo foi 88,4, sendo obtida por pelo menos um aluno da região Sul.
As demais notas máximas foram : 85,6 na região Nordeste, 83,6 na região Sudeste, 80,6 na região Norte, e 80,4 na região Centro-Oeste.
A nota mínima foi zero em todas as regiões, sem exceção.
Assim como os Gráficos 3.1 e 3.5, o Gráfico 3.9, apresentado a seguir, proporciona uma avaliação do desempenho de concluintes em relação ao componente de Conhecimento Específico com um histograma da distribuição das notas correspondentes.
Esta também é uma distribuição unimodal, e o grupo modal é o ( 30 ; 40 ], abaixo do grupo modal para a Formação Geral.
O coeficiente de assimetria da distribuição das notas do componente de Conhecimento Específico é positivo e próximo de zero ( 0,05 ).
Nota -se pelo histograma ( Gráfico 3.9 ) que esta é uma distribuição assimétrica à direita.
Nas regiões Nordeste ( 0,05 ), Sudeste ( 0,10 ) e Centro-Oeste ( 0,12 ) a assimetria mantém -se positiva.
Já nas regiões Norte e Sul, os coeficientes de assimetria são negativos, iguais a -0,05 e -0,21, respectivamente, evidenciando que a distribuição das notas dos estudantes dessas regiões, no componente de Conhecimento Específico, tem cauda mais pesada à esquerda.
Os Gráficos 3.10, 3.11 e 3.12 apresentam, respectivamente, uma comparação dos resultados em relação à Grande Região do país, à Categoria Administrativa e à Organização Acadêmica, agora levando em conta o desempenho de alunos no componente de Conhecimento Específico da prova.
Pelo Gráfico 3.10, observa -se que existe diferença estatisticamente significativa entre as médias das notas no componente de Conhecimento Específico, da região Sul em relação às demais regiões.
Sendo que esta região foi a que obteve média mais elevada ( 40,7 ), enquanto a região Centro-Oeste foi a região que obteve média mais baixa ( 35,2 ).
Quanto à Categoria Administrativa ( Gráfico 3.11 ), observa -se um comportamento semelhante ao da prova como um todo, ou seja, existe diferença estatisticamente significativa entre as médias das IES pública e privada, sendo que a das IES públicas é mais alta ( 41,2 ) do que a obtida pelos alunos das IES privada ( 36,5 ).
Quanto ao Gráfico 3.12, observa -se, mais uma vez, que existe diferença estatisticamente significativa ao nível de 95 \% entre as notas no componente de Conhecimento Específico dos diferentes tipos de Organização Acadêmica.
Sendo que a média dos Concluintes das Universidades ( 37,7 ) foi maior do que de Centros Universitários ( 35,7 ) e de Faculdades ( 36,3 ). 3.2 ANÁLISE DAS QUESTÕES OBJETIVAS 3.2.1 Componente de Formação Geral A Tabela 3.4 apresenta as estatísticas básicas relativas às oito questões objetivas do componente da prova que abrange a Formação Geral dos estudantes.
A média do Brasil foi 44,8.
A menor média foi encontrada na região Norte ( 43,2 ) e a maior na região Sul ( 46,7 ).
As demais médias foram 46,4 na região Nordeste, 44,8 na região Centro-Oeste e 44,1 na região Sudeste.
O desvio padrão do Brasil foi 20,9, sendo o maior desvio padrão encontrado nas regiões Norte e Centro-Oeste ( 21,2 ) e o menor na região Sul ( 20,6 ).
Os demais desvios foram : 21,0 na região Sudeste e 20,9 na região Nordeste.
A mediana ( 50,0 ) foi a mesma em todas as regiões, exceto na região Norte ( 37,5 ).
As mesmas notas máximas ( 100,0 ) e as mesmas notas mínimas ( 0,0 ), nas questões objetivas do componente de Formação Geral, foram obtidas em todas as regiões, sem exceção.
A Tabela 3.5 apresenta os índices de facilidade e discriminação ( ponto bisserial ) das questões objetivas do componente de Formação Geral.
Quanto ao índice de facilidade, foram usadas as seguintes cores para diferenciar o nível de dificuldade da questão : • Azul para as questões classificadas com índice muito fácil ( índice \&gt ; =0,86 ), verde para as questões classificadas com índice fácil ( 0,61 a 0,85 ), amarelo para as questões classificadas com médio ( 0,41 a 0,60 ), vermelho para as questões classificadas com difícil ( 0,16 a 0,40 ) e roxo para as questões classificadas com muito difícil ( \&lt ; =0,15 ).
Já quanto ao índice de discriminação, foram usadas as seguintes cores para qualificar a questão : • As questões classificadas com índice fraco receberam a cor roxa ( índice \&lt ; =0,19 ), as classificadas com médio receberam a cor vermelho ( 0,20 a 0,29 ), as classificadas com bom receberam a cor verde ( 0,30 a 0,39 ) e as classificadas com muito bom ( \&gt ; =0,40 ) receberam a cor azul.
As questões objetivas do componente de Formação Geral, segundo o índice de facilidade, foram assim avaliadas : das oito questões, nenhuma teve o índice de facilidade classificado como muito fácil, fácil ou muito difícil.
Seis questões foram tidas como médias, por terem índice de acertos situado na faixa entre 0,41 e 0,60 ( de 41,0 \% a 60,0 \% de acertos ).
Duas questões foram consideradas difíceis, situando -se no intervalo entre 0,16 e 0,40 do índice de facilidade, ou seja, houve entre 16,0 \% e 40,0 \% de acertos.
Como já comentado, para análise das questões objetivas relativas à Formação Geral segundo o poder de discriminação, utilizou -se, o índice de discriminação ponto bisserial.
Nesta análise as questões foram assim avaliadas : seis das oito questões apresentaram A Tabela 3.5 apresenta os índices de facilidade e discriminação ( ponto bisserial ) das questões objetivas do componente de Formação Geral.
Quanto ao índice de facilidade, foram usadas as seguintes cores para diferenciar o nível de dificuldade da questão :  Azul para as questões classificadas com índice muito fácil ( índice \&gt ; =0,86 ), verde para as questões classificadas com índice fácil ( 0,61 a 0,85 ), amarelo para as questões classificadas com médio ( 0,41 a 0,60 ), vermelho para as questões classificadas com difícil ( 0,16 a 0,40 ) e roxo para as questões classificadas com muito difícil ( \&lt ; =0,15 ).
Já quanto ao índice de discriminação, foram usadas as seguintes cores para qualificar a questão :  As questões classificadas com índice fraco receberam a cor roxa ( índice \&lt ; =0,19 ), as classificadas com médio receberam a cor vermelho ( 0,20 a 0,29 ), as classificadas com bom receberam a cor verde ( 0,30 a 0,39 ) e as classificadas com muito bom ( \&gt ; =0,40 ) receberam a cor azul.
As questões objetivas do componente de Formação Geral, segundo o índice de facilidade, foram assim avaliadas : das oito questões, nenhuma teve o índice de facilidade classificado como muito fácil, fácil ou muito difícil.
Seis questões foram tidas como médias, por terem índice de acertos situado na faixa entre 0,41 e 0,60 ( de 41,0 \% a 60,0 \% de acertos ).
Duas questões foram consideradas difíceis, situando -se no intervalo entre 0,16 e 0,40 do índice de facilidade, ou seja, houve entre 16,0 \% e 40,0 \% de acertos.
Como já comentado, para análise das questões objetivas relativas à Formação Geral segundo o poder de discriminação, utilizou -se, o índice de discriminação ponto bisserial.
Nesta análise as questões foram assim avaliadas : seis das oito questões apresentaram O Gráfico 3.13, para exemplificar, analisa o comportamento da questão de número 7 de Formação Geral.
Trata -se da terceira questão mais fácil e uma das que obteve o maior índice de discriminação dessa parte da prova.
Neste gráfico, \textbf{cada} uma das cinco curvas representa o percentual de respostas em determinada alternativa da questão, em função da nota dos estudantes nesta parte da prova ( Formação Geral/Múltipla Escolha ), antes de possíveis eliminações pelo critério do ponto bisserial.
A curva em verde corresponde à alternativa B, a correta para esta questão.
Assim, observa -se que entre os estudantes com notas mais baixas, nessa parte do exame, a situação mais frequente foi a escolha de uma das alternativas incorretas : a alternativa C ( em preto ) ou A ( em azul ).
Na medida em que a nota aumenta, indicando desempenho melhor nesta parte da prova, aumenta concomitantemente a proporção de estudantes que índices acima de 0,40 e, assim, foram classificadas com índice muito bom para esse grupo de alunos ; uma questão teve índice de discriminação bom, entre 0,30 e 0,39.
Foi classificada com índice médio apenas uma das questões.
O índice de facilidade variou de 0,28 a 0,59, e o de discriminação, de 0,29 a 0,48.
As questões com índices de discriminação muito bom, de números 1, 2, 3, 4, 6 e 7, figuraram entre as mais fáceis desse conjunto : todas classificadas na categoria médio.
Em particular, as questões 6 e 7 foram as que apresentaram maior poder discriminatório, com índice 0,48, e foram também a segunda e a terceira mais fáceis, com uma proporção de 0,56 e 0,49 de acertos, respectivamente.
A questão de número 5 apresentou índice de facilidade 0,28, ou seja, 28 \% dos estudantes conseguiu resolvê -la, dentro do universo de participantes.
Além disso, seu índice de discriminação foi médio.
A questão 8 obteve índice de discriminação bom, 0,37, e seu índice de facilidade foi difícil.
O Gráfico 3.13, para exemplificar, analisa o comportamento da questão de número 7 de Formação Geral.
Trata -se da terceira questão mais fácil e uma das que obteve o maior índice de discriminação dessa parte da prova.
Neste gráfico, \textbf{cada} uma das cinco curvas representa o percentual de respostas em determinada alternativa da questão, em função da nota dos estudantes nesta parte da prova ( Formação Geral/Múltipla Escolha ), antes de possíveis eliminações pelo critério do ponto bisserial.
A curva em verde corresponde à alternativa B, a correta para esta questão.
Assim, observa -se que entre os estudantes com notas mais baixas, nessa parte do exame, a situação mais frequente foi a escolha de uma das alternativas incorretas : a alternativa C ( em preto ) ou A ( em azul ).
Na medida em que a nota aumenta, indicando desempenho melhor nesta parte da prova, aumenta concomitantemente a proporção de estudantes que selecionaram a alternativa correta B, atingindo 100 \% para os alunos com 8 acertos.
Essa análise permite verificar como a questão discriminou os grupos de desempenho, justificando o alto índice obtido na questão.
Os gráficos relativos às demais questões de Formação Geral constam do Anexo I. 3.2.2 Componente de Conhecimento Específico A Tabela 3.6 apresenta as estatísticas básicas em relação às questões objetivas do componente de Conhecimento Específico da prova, por Grande Região.
A média do Brasil deste componente foi de 41,1.
A menor média foi observada na região Norte ( 39,4 ) e a maior na região Sul ( 45,2 ).
O desvio padrão de todo o Brasil foi 16,2, sendo o menor desvio padrão encontrado na região Norte ( 15,1 ) e o maior na região Sul ( 16,4 ).
A mediana de todo o Brasil foi 41,7 e a mesma em quase todas as regiões, exceto nas regiões Sul ( 45,8 ) e Centro-Oeste ( 37,5 ).
A nota máxima da prova ( 91,7 ) foi obtida, nas questões objetivas do componente de Conhecimento Específico, por pelo menos um aluno da região Sul.
A nota máxima nas demais regiões foi : 87,5 no Nordeste, Sudeste e Centro-Oeste e 83,3 na região Norte.
A nota mínima foi zero em todas as regiões, sem exceção.
A Tabela 3.7 apresenta os índices de facilidade e discriminação ( ponto bisserial ) das questões objetivas do componente de Conhecimento Específico.
Para facilitar a diferenciação das questões usaram -se as mesmas cores da Tabela 3.5 para as diferentes classificações dos índices de facilidade e de discriminação.
Dentre as questões objetivas da parte da prova relativa ao componente de Conhecimento Específico, nenhuma foi anulada pela Comissão.
Desse modo, a classificação quanto ao índice de facilidade foi estabelecida com base em todas as 27 questões.
A partir dos índices obtidos, pode -se concluir que a maioria das questões objetivas da prova foi considerada pelo menos difícil : das 27 questões, quinze foram classificadas como difíceis e duas como muito difíceis.
Não houve questão classificada como muito fácil, ao passo que duas foram tidas como fáceis, na faixa de 0,61 a 0,85 do índice de facilidade, e outras oito consideradas médias, entre 0,41 e 0,60.
Já quanto aos índices de discriminação das questões objetivas do componente de Conhecimento Específico da prova, tem -se como resultado a seguinte classificação : treze das 27 questões foram consideradas como boas, enquanto seis delas tiveram índice de discriminação muito bom.
Assim, para 19 das 27 questões, os índices de discriminação foram bons ou muito bons.
Dentre as demais, cinco delas foram classificadas como médias e outras três como fracas, sendo oito, por conseguinte, a quantidade de questões nos dois patamares mais baixos de discriminação.
Constata -se, assim, que a prova – no que se refere ao componente de Conhecimento Específico – possuía capacidade boa de discriminar entre aqueles que dominam ou não o conteúdo.
Dentre as questões que alcançaram os maiores índices de discriminação, as de números 20, 24, 27, 29, 33 e 34 classificadas com índice muito bom, a partir de 0,40 do índice, duas delas ( questões 20 e 29 ) foram classificadas na categoria fácil, três ( 27, 33 e 34 ) como médias e uma ( 24 ) como difícil, quanto ao índice de facilidade.
A questão de número 28 foi a mais difícil dentre as 27 questões específicas, com baixo índice de facilidade, apenas 10,0 \% de acertos.
Essa questão apresentou poder discriminatório igualmente baixo, -0,08, o que comprova ter sido esta a mais difícil para os estudantes.
Destaca -se, também, a questão 35, com índice de facilidade 0,16, o que, em termos percentuais, corresponde a 16,0 \% de estudantes que responderam acertadamente, obtendo, ainda, 0,15 de índice de discriminação.
Tais questões foram, portanto, duas das mais difíceis da prova.
As questões 28 e 35 foram eliminadas do cômputo da nota final pelo critério do ponto bisserial.
Além destas duas, a questão 31, com índice fraco de discriminação, também não foi computada.
A título de exemplo das análises do comportamento das questões objetivas, o Gráfico 3.14 analisa a questão 29 do componente de Conhecimento Específico.
Esta foi a questão mais fácil da prova, apresentando índice de facilidade 0,74, ou seja, 74,0 \% dos estudantes assinalaram acertadamente a opção D, correspondente ao gabarito.
Seu índice de discriminação foi igual a 0,46, classificado como muito bom.
Neste gráfico, \textbf{cada} uma das cinco curvas representa o percentual de respostas em determinada alternativa da questão 29, em função da nota dos estudantes nesta parte da prova, antes de possíveis eliminações de questões pelo critério do ponto bisserial.
A alternativa correta D, representada no gráfico pela curva em roxo, foi escolhida em maiores proporções pelos alunos com desempenho melhor nesta parte da prova.
Já as alternativas incorretas, também denominadas distratores, foram selecionadas principalmente por Aqueles com nota zero, na sua quase totalidade deixaram esta questão em branco ou marcaram mais de uma alternativa, comportamento considerado inválido.
A proporção de alunos que selecionou a resposta correta D aumenta gradativamente, chegando a atingir 100 \% a partir de 22 acertos, enquanto a proporção dos que escolheram alternativas incorretas decai, a partir da primeira nota não nula, como função da nota nesta parte da prova.
Os gráficos relativos às demais questões do componente de Conhecimento Específico constam do Anexo I. 3.3 ANÁLISE DAS QUESTÕES DISCURSIVAS 3.3.1 Componente de Formação Geral As análises dos resultados de desempenho dos estudantes de Comunicação Social - Publicidade e Propaganda nas duas questões discursivas relativas à Formação Geral encontram -se na Tabela 3.8 e no Gráfico 3.15.
Na tabela 3.8 observa -se que as notas médias foram mais baixas nesse conjunto de questões do que no das objetivas.
Os estudantes de todo o Brasil obtiveram, em Formação Geral, média 44,8 nas questões objetivas e 34,5 nas questões discursivas.
A mediana também confirma o pior desempenho dos alunos de todo o Brasil nas questões discursivas do componente de Formação Geral.
Enquanto essa estatística foi de 37,5 para questões discursivas, para as questões objetivas essa estatística foi de 50,0.
Pode -se, também, notar um aumento do desvio padrão de 20,9, nas questões objetivas do componente de Formação Geral dos alunos de todo o Brasil, para 22,7 nas questões discursivas do mesmo componente.
Como já comentado, a mediana de todo o Brasil, neste componente, foi 37,5, sendo a maior mediana encontrada na região Sudeste ( 42,5 ) e a menor na região Centro-Oeste ( 30,0 ).
A nota máxima ( 92,5 ) foi obtida, nas questões discursivas do componente de Formação Geral, por pelo menos um aluno na região Sudeste.
A nota máxima das demais regiões foi 82,5 das Norte e Centro-Oeste, 87,5 da Nordeste e 90,0 da Sul.
A nota mínima zero foi obtida em todas as regiões, sem exceção.
O Gráfico 3.15 representa a distribuição das notas nas questões discursivas no componente de Formação Geral.
A moda desta distribuição ocorre no primeiro intervalo, [ 0 ; 10 ], devido à grande quantidade de notas zero e à alta frequência de alunos que deixaram este tipo de questão em branco.
Desconsiderando esse intervalo, a distribuição possui outro intervalo modal, ( 40 ; 50 ].
A distribuição possui assimetria à esquerda, com coeficiente de assimetria -0,24.
Em todas as regiões o coeficiente de assimetria é negativo, exceto na região Centro-Oeste ( 0,08 ).
O maior coeficiente foi obtido na região Centro-Oeste ( 0,08 ) e o menor na região Sul ( -0,40 ).
Na sequência, os resultados verificados para \textbf{cada} uma das questões discursivas de Formação Geral serão apresentados, estabelecendo -se relações com os conteúdos abordados em \textbf{cada} uma delas.
Os comentários da Banca de docentes corretores a respeito do observado na correção das respostas dos estudantes, suas impressões e conclusões serão apresentados junto à análise de \textbf{cada} questão.
Cumpre esclarecer que, tendo em vista que as questões discursivas de Formação Geral são padronizadas, ou seja, constam de todas as provas, os comentários da Banca são os mesmos para todas as carreiras acadêmicas, sendo direcionados a todos os estudantes que participaram do ENADE/2012.
A seguir, serão analisados os desempenhos da Área de Comunicação Social - Publicidade e Propaganda nas duas questões discursivas de Formação Geral do ENADE/2012, comparando os resultados obtidos com comentários para \textbf{cada} questão. 3.3.1.1 Análise da Questão Discursiva 1 do Componente de Formação Geral Os dados de Comunicação Social – Publicidade e Propaganda, obtidos a partir das respostas à questão 1, encontram -se na Tabela 3.9 e no Gráfico 3.16.
Nessa questão – de pior desempenho dentre as duas de Formação Geral – os alunos de todo Brasil tiveram média, 30,3.
A maior média para a questão 1 foi obtida na região Sul ( 32,7 ), e a menor, na região Centro-Oeste ( 25,6 ).
Quanto à variabilidade das notas, o desvio padrão de todo o Brasil foi 25,1.
O menor desvio padrão foi obtido nas regiões Sudeste e Sul ( 24,8 ) e o maior foi obtido nas regiões Norte e Nordeste ( 26,2 ).
As medianas e as notas máximas da questão discursiva 1 foram as mesmas para quase todas as regiões do Brasil, respectivamente, 30,0 e 95,0, exceto na região Centro-Oeste, a mediana ( 20,0 ) e na região Norte, a nota máxima ( 90,0 ).
Além disso, a mesma nota mínima ( 0,0 ) foi obtida em todas as regiões do país, sem exceção.
O Gráfico 3.16 mostra a distribuição das notas na questão discursiva 1 do componente de Formação Geral.
Observa -se que a maior frequência corresponde aos alunos que deixaram a questão em branco, correspondendo à moda.
Considerando apenas os estudantes que tentaram resolver a questão, a moda ocorre no intervalo ( 30 ; 40 ].
Também observa -se grande concentração de notas nos intervalos ( 20 ; 30 ] e ( 40 ; 50 ].
O coeficiente de assimetria para todo o Brasil foi positivo ( 0,25 ), indicando uma assimetria à direita.
Em todas as regiões o coeficiente de assimetria também foi positivo ( entre 0,15 para a região Sul e 0,46 para a região Centro-Oeste ). 3.3.1.2 Comentários sobre a correção das respostas à Questão Discursiva 1 A questão aborda um tema bastante polêmico o qual vem sendo amplamente discutido, a conciliação entre o desenvolvimento econômico e a preservação dos recursos naturais.
Busca -se observar a relação entre o estímulo governamental à compra de veículos novos e os impactos ambientais decorrentes desta política pública.
Sua compreensão implica no conhecimento de uma temática atual e na reflexão critica a respeito das decisões governamentais em sua correlação com o cenário econômico internacional e os problemas e soluções de \textbf{ordem} ambiental.
Embora o comando da questão solicitasse explicitamente que a resposta fosse fornecida em um texto único ( “ redija um texto dissertativo sobre... ” ), foi significativamente alta a proporção de respondentes que atendeu a \textbf{cada} item ( a, b e c ) do enunciado de maneira isolada, como partes estanques da mesma resposta.
No que se segue encontra -se uma avaliação da correção de \textbf{cada} um dos três itens. a - “ conceito de desenvolvimento sustentável ” A variedade de respostas a esse item confirmou a importância do tema para avaliar a formação geral dos estudantes e cursos, já que ( i ) é atual ; ( ii ) foi abordado nos meios de comunicação ( especialmente no ano de 2012, em que o Brasil sediou a “ Rio+20 ” ) ; ( iii ) é relevante não apenas no Brasil como em todos os principais fóruns internacionais ; e ( iv ) é comum a diversas áreas de conhecimento.
Outros dois aspectos positivos merecem destaque : ( i ) a grande incidência de respostas que apontaram a relação entre gerações como aspecto essencial da ideia de sustentabilidade socioambiental – no mais, as respostas apresentaram pelo menos a tentativa de indicar uma definição de desenvolvimento sustentável ou a abordagem positiva dessa mesma noção, evidenciado compreensão do enunciado ; e ( ii ) o exercício crítico que a questão propunha, no sentido de exigir dos respondentes a compreensão das contradições de políticas públicas e a apresentação de propostas de soluções para conciliar o incentivo do crescimento econômico e a promoção da sustentabilidade socioambiental, a partir de uma perspectiva transformadora da realidade.
Apesar dos aspectos positivos apontados no parágrafo anterior, a maioria dos estudantes não alcançou o máximo da pontuação.
De forma geral, trataram o desenvolvimento sustentável como sinônimo de não agressão ambiental.
Alguns confundiram desenvolvimento sustentável com desenvolvimento econômico, e outros faziam referência apenas à preocupação com a conservação dos recursos mencionando a preocupação com gerações futuras.
Deve -se, ainda, destacar o significativo número de alunos que apresentou uma definição aparentemente copiada do enunciado de outra questão da prova, como por exemplo, a questão objetiva 31 da prova de Administração. b - “ conflito entre o estímulo à compra de veículos automotores e a promoção da sustentabilidade ” Esse foi o item identificado como aquele no qual as respostas dos estudantes avaliados mais se aproximaram da pontuação máxima prevista no padrão de resposta.
Ou seja, a ampla maioria dos respondentes foi bem-sucedida em apontar a relação entre o aumento da quantidade de veículos automotores em circulação e os danos socioambientais que isso acarreta no cenário urbano, devido ao aumento da emissão de gases poluentes.
No item b, as respostas corretas se concentraram em críticas ao incentivo oferecido pelo governo, referências ao problema dos congestionamentos e ao aumento da poluição.
Contudo, a maioria das respostas não explicitava a correlação entre o aumento da circulação de veículos automotores e o agravamento dos problemas ambientais.
Um número expressivo de alunos apresentou uma resposta em que era dado destaque à política de redução do IPI como um mecanismo para renovação da frota de automóveis ou para a aquisição do transporte próprio e da consequente redução da dependência ao transporte coletivo.
Também foram notados alguns tipos de padrões recorrentes nas tentativas de resposta a este item.
Dois desses padrões se destacaram : ( i ) as declarações genéricas de condenação do “ consumismo ” de uma “ nova classe média ” ou da busca pelo “ status ” que a posse de um automóvel particular proporcionaria, como se fosse possível atribuir a essa \textbf{ordem} de coisas a causa principal do aumento dos congestionamentos nas grandes cidades ; e ( ii ) a repetição da avaliação, presente no enunciado da questão, que aponta de forma superficial a contradição entre o estímulo à demanda por automóveis particulares e a promoção da sustentabilidade. c - “ ações de fomento ao transporte urbano sustentável no Brasil ” Este foi o item que ofereceu maior facilidade de aplicação do padrão de resposta, mas foram poucos os casos de pontuação integral.
Dentre as soluções mais destacadas pelos alunos apareceram : ( i ) o incentivo ao uso de energias alternativas ; e ( ii ) o estímulo ao transporte coletivo ( metrô, trem e ônibus ), além do uso de bicicletas.
Importante destacar o fato de que muitos apresentavam respostas com referências à necessidade de incentivo ao transporte urbano, nas quais este termo era utilizado como sinônimo de transporte coletivo.
A avaliação da questão 1 trouxe à tona problemas estruturais da educação básica, que se refletem no ensino superior.
Foi muito grande o número de provas em que o aluno demonstrou não compreender o enunciado da questão, apresentando respostas que fugiam ao tema proposto, mesmo quando bem redigidas.
Uma resposta recorrente fazia referência às consequências da ampla oferta de crédito à população e aos problemas relacionados à inadimplência, situação que, embora, verdadeira, não \textbf{estava} relacionada à questão proposta.
Houve, ainda, os casos em que as respostas apresentadas não só se afastaram do tema como não reuniam qualquer tipo de argumentação coerente.
E, por fim, destaca -se a existência de respostas incoerentes e desconexas.
Outro aspecto geral diz respeito ao grande número de provas em branco ou em que o aluno apresentou a resposta a outro item no espaço da questão 1.
Considerando -se o objetivo geral deste exame, avaliar a qualidade dos cursos de graduação oferecidos no país e contribuir para a melhoria da qualidade do ensino superior, destaca -se a imensa responsabilidade do processo que envolve a elaboração deste diagnóstico.
Nesse sentido, é urgente alertar para a ocorrência de um número significativo de respostas que se afastam de um padrão mediano de qualidade esperado de um aluno concluinte em um curso de graduação, especialmente em se considerando o fato de que a questão proposta faz referência a tema intensamente explorado pela mídia. 3.3.1.3 Análise da Questão Discursiva 2 do Componente de Formação Geral A Tabela 3.10 mostra que o desempenho dos estudantes na questão 2 ( média 38,6 ) foi superior ao obtido na questão de número 1 ( média 30,3 ).
A região Norte foi aquela onde a média, nessa questão, foi maior ( 43,6 ), e a de menor média foi a região Centro-Oeste ( 32,3 ).
Quanto à variabilidade das notas, o desvio padrão de todo o Brasil foi 26,5, superior ao obtido na questão de número 1 ( 25,1 ).
O maior desvio nessa questão foi obtido na região Centro-Oeste ( 28,0 ), enquanto o menor foi obtido na região Sul ( 25,9 ).
A mediana comprova a diferença de desempenho entre essas duas regiões : enquanto a menor mediana ocorreu na região Centro-Oeste ( 40,0 ), a maior ocorreu nas regiões Norte e Sul ( 50,0 ).
A nota máxima ( 95,0 ) foi obtida em quase todas as regiões do Brasil, exceto na região Norte ( 90,0 ).
A nota mínima ( 0,0 ) foi a mesma em todas as regiões, sem exceção.
O Gráfico 3.17 mostra a distribuição das notas na questão discursiva 2 do componente de Formação Geral.
Neste gráfico, nota -se uma grande quantidade de alunos que deixaram a questão 2 em branco ( em \textbf{torno} de 24 \% ), sendo esta a situação modal.
Considerando apenas os estudantes que tentaram resolver a questão, a moda ocorre no intervalo ( 40 ; 50 ].
O coeficiente de assimetria das notas nesta questão é negativo ( -0,32 ), indicando que a distribuição é assimétrica à esquerda.
O mesmo ocorre para quase todas as regiões ( entre -0,53 para a região Norte e -0,24 para a região Nordeste ), com exceção da região Centro-Oeste onde o coeficiente é positivo ( 0,02 ), ainda que bem perto de zero. 3.3.1.4 Comentários sobre a correção das respostas à Questão Discursiva 2 O conteúdo da questão era bastante adequado às Diretrizes do INEP para o ENADE, além de atual e de abrangente.
Porém, exigia -se a articulação de muitas variáveis.
Buscando aferir, a partir da apresentação de uma charge e de um texto introdutório, as competências e habilidades relacionadas à leitura, à interpretação, à análise crítica, à capacidade argumentativa, à elaboração de sínteses e de avaliação prospectiva sobre tema da atualidade, a questão tratava da violência e suas repercussões no contexto atual.
O comando da questão solicitava, a partir da definição de violência da OMS e de duas imagens, que o aluno elaborasse um texto, de caráter dissertativo, sobre a “ violência na atualidade ”.
Embora tal solicitação fosse explícita ( “ redija um texto dissertativo sobre... ” ), como na questão 1, foi alta a proporção de alunos que respondeu a \textbf{cada} item do enunciado isoladamente, como partes independentes.
Verificaram -se poucas respostas em branco, como também foram poucas as respostas contendo protestos contra o ENADE, muito frequentes em anos anteriores, o que significa um avanço para a afirmação do ENADE como um exame aceito socialmente. a -"tecnologia e violência"Tanto as charges quanto o texto apresentado no enunciado induziram os alunos a focarem suas respostas nos meios de comunicação.
Além disso, algumas respostas estabeleciam a tecnologia, em si, como causa da violência tanto nas escolas ( em particular ), quanto na sociedade ( no geral ).
Muitos respondentes tiveram dificuldade para fazer a relação entre tecnologia, violência e a escola, de modo adequado, centrando -se na violência nas escolas.
Além disso, a grande maioria das respostas girava em \textbf{torno} da exposição da violência através da TV, internet e jogos digitais, sendo que predominaram as abordagens com ênfase em programas e filmes veiculados pela televisão.
Poucos foram os estudantes que estabeleceram relação entre o contexto socioeconômico e a violência. b -"causas e consequências da violência na escola"Embora a solicitação desse item fosse objetiva e atendida por grande parte dos alunos, inúmeras vezes, observou -se confusão entre o significado de causa e de consequência.
Em outras inúmeras vezes, as respostas apontavam apenas para uma das classes, causas ou consequências.
Foram frequentes as respostas que atribuíam o aumento da violência à ineficiência do poder público, apontando a falta de investimentos, o baixo salário dos professores, a ausência de profissionais qualificados como causas.
Foram frequentes também os que apontavam a violência doméstica, o bullying, e a transferência do imaginário dos filmes e jogos para a realidade, como causas da violência. c -"proposta de solução para o problema da violência na escola"Sem dúvida, há uma grande diversidade de abordagens possíveis para apresentação de uma proposta de solução para o problema da violência, o que o padrão de resposta buscou contemplar.
No entanto, houve grande número de respostas baseadas em opiniões pessoais dos estudantes e vinculadas ao senso comum.
As respostas giraram em \textbf{torno} da necessidade de mais investimentos na educação por parte do poder público e seu rebatimento no maior aparelhamento das escolas, salário dos professores etc.
Além disso, foram muitas as respostas que trataram da necessidade de maior conscientização dos pais, alunos e professores com a realização de um trabalho conjunto que envolvesse todos os segmentos da sociedade.
De um modo geral, os estudantes demonstraram dificuldades em relação ao nível de abrangência exigido, falta de conteúdo, inabilidade na elaboração de texto escrito.
Da grande quantidade de provas corrigidas, o que ficou evidente, foi a dificuldade de estabelecer as conexões requeridas na questão 2.
A violência nas escolas foi bastante associada ao reflexo do meio onde a criança e o adolescente mantém os seus vínculos afetivos.
Surgiram muitas respostas revelando críticas à intolerância no núcleo familiar, à jornada de trabalho dos pais, ao distanciamento da mãe nas atividades domésticas etc.
Em boa parte das respostas, os estudantes lançaram a causa da violência escolar exclusivamente na ausência de participação familiar na vida dos filhos, às vezes criando uma cisão com a proposta da questão, já que abandonavam totalmente a temática da tecnologia em suas reflexões.
Além da ausência dos pais no cotidiano da escola, divórcios e atos de violência contra crianças foram considerados fatores que delegariam para a escola novas funções.
Também foram mencionadas a má formação de professores e a insatisfação com o ambiente escolar.
Em inúmeros casos, ficou evidente o anseio de se fazer da escola um lugar de resgate tanto da criança e do adolescente quanto da família, compreendendo -se a escola em sentido social e assistencial.
Segundo este amplo grupo de concluintes, a escola deveria oferecer -se como centro de cultura e lazer, como espaço para o acompanhamento psicológico e serviço social às comunidades diretamente envolvidas.
Ressalte -se o constante aparecimento das sugestões de horário integral, melhores salários para professores e iniciativas para melhoria na formação de professores.
As propostas de solução para acabar com a violência na escola se detiveram em aspectos muito gerais e poucos alunos detalharam uma proposta de intervenção, apresentando ações concretas.
Dentre as propostas de solução surgiram : educação familiar, ensino religioso, maior segurança nas escolas, com utilização de aparelhos para detecção de armas, e o retorno do papel do professor como um profissional respeitado pelos poderes municipais e estaduais.
Um percentual significativo de provas apresentou resposta simplista, com ênfase em determinados pontos como : religião ( sem demonstração do conceito de laicidade ), políticas sociais ( sem demonstração do domínio de conceitos que envolvem Estado, Governo e Politicas Públicas ) etc.
Foi grande a incidência de respostas que se referiam às leis, à vigilância e maior controle dos pais e professores ; no entanto, não demonstram entender a noção de legalidade para os diferentes segmentos sociais e como os mesmos se encontram submetidos a diferentes formas de poder.
Alguns alunos, ainda que raros, trataram a questão com a profundidade esperada de concluintes de cursos universitários.
Por outro lado, foram vários os textos que apresentaram uma argumentação simplória e desconexa, com forte visão paternalista sobre o papel do governo e do professor, como se tais instâncias pudessem intervir na educação familiar.
Enquanto alguns conseguiram desenvolver o tema proposto com elegância e desenvoltura, encadeando o tema de forma harmônica, com excelente nível de linguagem escrita e correção gramatical, outros sequer conseguiram conectar as ideias e construir um texto razoável, utilizando, inclusive, um linguajar mais propício à linguagem falada do que à escrita.
Cabe também ressaltar a quantidade de respostas carregadas de conteúdo com juízo de valor discriminatório ou extremamente conservador, do tipo : “ mulher não deve trabalhar fora ” ; “ amar e temer a Deus é a solução ” ; obrigatoriedade do ensino religioso e de uma doutrina religiosa ; imposição de uma educação do tipo militar nas escolas ; censura nos meios de comunicação etc.
Outras soluções apontadas para resolver o problema da violência nas escolas \textbf{estavam} associadas à, por exemplo : redução da maioridade penal, defesa da pena de morte, milícias armadas na porta das escolas, punições como o retorno da palmatória.
O Estatuto da Criança e do Adolescente foi uma peça legal criticada, e recebeu, em diversas respostas, o título de responsável pelo sentimento de impunidade de adolescentes infratores.
Contudo, algumas dessas proposições, desde que bem contextualizadas e formuladas com boa argumentação, foram pontuadas de forma a respeitar o direito democrático de pluralidade opinativa.
É imperioso ressaltar que a falta de inteligibilidade segue representando um aspecto merecedor de providências na formação universitária, uma vez que exprime uma habilidade importante para qualquer profissional de nível superior.
Além de terem sido poucas respostas boas e claras, foi possível identificar a deficiência de argumentação, além dos inúmeros problemas gramaticais ( ortografia, concordância, pontuação, regência, etc. ).
Salienta -se, nesse sentido, o uso indiscriminado de letras maiúsculas e minúsculas numa mesma frase, sem obedecer ao critério formal da língua portuguesa. 3.3.2 Componente de Conhecimento Específico Na parte da prova relativa às questões discursivas no componente de Conhecimento Específico ( Tabela 3.11 ), observa -se que a média foi mais baixa do que para as questões discursivas do componente de Formação Geral.
Enquanto no componente de Formação Geral a média para estudantes de Comunicação Social - Publicidade e Propaganda de todo o Brasil foi 34,5, na parte de Conhecimento Específico a média foi 12,8.
A maior média deste componente foi obtida pelos estudantes da região Sul ( 15,0 ), e a menor, pelos da região Centro-Oeste ( 9,8 ).
Quanto à variabilidade das notas, o desvio padrão de todo o Brasil foi 15,6.
O maior desvio padrão foi encontrado na região Sul ( 17,1 ), e o menor, na região Centro-Oeste ( 14,5 ).
A maior nota máxima, nas questões discursivas do componente de Conhecimento Específico, foi obtida na região Sudeste ( 86,7 ) enquanto a menor nota máxima foi encontrada na região Norte ( 65,0 ).
Além disso, a nota mínima ( 0,0 ) foi obtida por alunos de todas as regiões do Brasil sem exceção.
A mediana do Brasil como um todo foi 6,7, sendo a maior obtida na região Norte ( 11,7 ) e a menor obtida na região Centro-Oeste ( 0,0 ).
Ou seja, pelo menos metade dos estudantes da região Centro-Oeste tirou zero nesta questão.
O Gráfico 3.18 representa a distribuição das notas nas questões discursivas no componente de Conhecimento Específico.
Essa distribuição tem moda no intervalo de [ 0 ; 10 ].
Nota -se que a partir desse intervalo há uma outra concentração de notas no intervalo ( 10 ; 20 ], também sendo relevante não terem sido registradas notas acima de 90,0.
O coeficiente de assimetria é positivo ( 1,33 ), indicando uma assimetria à direita.
Em todas as regiões do Brasil o coeficiente de assimetria também é positivo, sendo o maior o da região Centro-Oeste ( 1,75 ) e o menor o da região Norte ( 0,91 ). 3.3.2.1 Análise da Questão Discursiva 3 do Componente de Conhecimento Específico Na questão 3, cujos resultados aferidos encontram -se descritos na Tabela 3.12, a média dos estudantes de todo o Brasil foi 17,9.
A menor média nessa questão foi obtida pelos alunos da região Centro-Oeste ( 14,0 ), enquanto a maior média foi obtida na região Norte ( 21,9 ).
Quanto à variabilidade das notas, o desvio padrão de todo o Brasil foi 24,0.
O maior desvio padrão foi obtido na região Sul ( 25,5 ), enquanto o menor foi obtido na região Centro-Oeste ( 22,3 ).
A mesma nota máxima, 100,0 pontos, mediana ( 0,0 ) e nota mínima ( 0,0 ) foram obtidas em todas as regiões do Brasil, com exceção da região Norte, cuja mediana foi 15,0.
Ou seja, na maioria das regiões do país pelo menos metade dos estudantes tirou zero nesta questão.
O Gráfico 3.19 mostra a distribuição das notas na questão discursiva 3, do componente de Conhecimento Específico.
Essa distribuição é bimodal, com moda principal nas questões em branco ( também com nota zero ) e uma segunda moda no intervalo ( 20 ; 30 ].
O coeficiente de assimetria é positivo ( 1,12 ).
Em todas as regiões o coeficiente de assimetria também é positivo ( entre 0,72 para a região Norte e 1,45 para a região Centro-Oeste ). 3.3.2.2 Comentários sobre as respostas à Questão Discursiva 3 O enunciado e a linguagem utilizada eram claros, não deixando margem a interpretações diferentes pelo candidato.
O conteúdo abordado pareceu um pouco extenso e profundo, levando em consideração a relação entre o que é pedido na questão, o espaço para resposta e o tempo de prova.
A primeira observação a ser feita em relação à questão 3 diz respeito à incapacidade do desenvolvimento da discussão teórica por parte dos alunos.
Tratando de modo geral, a aplicação das teorias da comunicação na dissertação acerca da televisão como veículo de cultura \textbf{esteve} quase sempre restrita à argumentação rasa, superficial, baseada mais no senso comum do que no conhecimento do assunto.
As respostas trouxeram, em sua maioria, apenas críticas individuais ao conteúdo da grade de programação das emissoras, repudiando o excesso de banalidade, futilidade, violência e “ sexo explícito ” na TV aberta e dando como alternativa os canais da TV fechada, que conteriam conteúdo considerado cultural.
Quando citados autores e teorias ( a maioria descreveu a Indústria Cultural de Adorno e Horkheimer querendo o público como mero consumidor ), era apenas para justificar a crítica feita de modo solto, sem aprofundar a discussão ou embasar um ponto de vista.
A segunda observação dá conta de uma parcela dos alunos que limitou a discussão a uma relação entre programação televisiva violenta e público infantil, repudiando a maneira despreocupada com que as emissoras \textbf{estariam} veiculando seus conteúdos às crianças.
Uma terceira observação é que a grande quantidade de alunos entendeu a tirinha como uma má interpretação da Mafalda quanto ao termo veículo.
A personagem teria confundido veículo de cultura com veículo automotor, fazendo assim seu comentário final, que segundo os alunos não deixava de ser ‘ verdadeiro ’.
Houve ainda um número muito significativo de respostas que começavam discordando da tirinha, pois consideravam a TV um veículo de cultura, mas ao final deixavam bem claro a concordância com a crítica da personagem ‘ pois a TV hoje em dia não exibiria conteúdo cultural relevante ’.
Nesse ponto, percebemos nitidamente a dificuldade dos alunos em organizar um raciocínio crítico somado ao seu ponto de vista quando exigidos três níveis de informação ( os subitens da questão – TV como veículo de cultura ; autores e teorias da comunicação ; ponto de vista do aluno ).
Vale ressaltar a grande dificuldade gramatical observada nas respostas apresentadas.
O contexto de aplicação da prova poderia explicar ( ainda que não concordemos com essa justificativa ) a grande quantidade de questões em branco e, em algum grau, o descompromisso com as respostas, mas nunca o desconhecimento da língua portuguesa da maneira como encontrado nas provas. 3.3.2.3 Análise da Questão Discursiva 4 do Componente de Conhecimento Específico A Tabela 3.13 contém as informações relativas à questão 4 do conjunto de questões do componente de Conhecimento Específico.
O desempenho dos estudantes de todo o Brasil nesta questão foi inferior ao desempenho na questão de número 3.
A média geral do Brasil foi 11,8, sendo a menor média registrada na região Centro-Oeste ( 8,2 ) e a maior na região Sul ( 14,9 ).
A mesma nota máxima ( 100,0 ), mediana ( 0,0 ) e nota mínima ( 0,0 ) foram obtidas em todas as regiões do Brasil, sem exceção.
Ou seja, pelo menos metade dos estudantes tirou zero nesta questão.
O Gráfico 3.20, representa a distribuição de notas na questão discursiva 4, no componente de Conhecimento Específico.
Como nas outras questões discursivas, o número de estudantes que deixou a questão em branco foi muito grande, correspondendo à moda da distribuição.
Dentre as notas dos que resolveram a questão 4, o intervalo modal foi das notas zero.
O coeficiente de assimetria é positivo ( 2,03 ), indicando uma concentração do lado esquerdo da distribuição e uma cauda longa do lado direito deste intervalo.
Tanto na distribuição como um todo quanto na análise dos resultados dos alunos por região, os coeficientes de assimetria são sempre positivos e com valor alto.
A região Sul ( 1,66 ) e a região Centro-Oeste ( 2,73 ) obtiveram o menor e o maior coeficiente de assimetria, respectivamente. 3.3.2.4 Comentários sobre as respostas à Questão Discursiva 4 A estrutura da questão levou muitos alunos, que desconheciam as teorias da Galáxia de Gutenberg e da Galáxia da Internet, a fazerem uma simples interpretação do texto, sem apresentar definições teóricas e muito menos a comparação crítica esperada.
Isso tornou a correção mais complexa : os alunos não \textbf{estavam} completamente errados em suas afirmações, mas não atendiam ao solicitado pela questão.
Os alunos que não conheciam as teorias tentavam, a partir do texto de Castells, formular uma resposta.
Com conteúdo mais raso, o aluno repetia as palavras e termos do texto de referência.
No entanto, como o padrão de resposta pedia que o aluno apresentasse os benefícios e a influência da internet na sociedade, e ele chegava a isso, mesmo sem a comparação aprofundada entre as duas teorias, acabava alcançando pontos na questão. 3.3.2.5 Análise da Questão Discursiva 5 do Componente de Conhecimento Específico A Tabela 3.14 contém as informações relativas à questão 5 do conjunto do componente de Conhecimento Específico.
O desempenho dos estudantes nessa questão foi inferior ao das questões 3 e 4.
A nota média dos estudantes de todo o Brasil foi 8,7.
A maior média foi registrada na região Sul ( 10,3 ), enquanto a menor média foi registrada na região Centro-Oeste ( 7,1 ).
Quanto à variabilidade das notas, o desvio padrão dos alunos do Brasil, como um todo, foi 14,4.
Enquanto o maior desvio foi encontrado na região Sul ( 15,5 ), o menor foi encontrado na região Sudeste ( 13,9 ).
A nota máxima ( 100,0 ) foi alcançada por pelo menos um estudante na região Sudeste, enquanto que nas demais regiões as notas máximas foram : 95,0 na região Nordeste, 90,0 na região Norte, 85,0 na região Sul e 75,0 na região Centro-Oeste.
A mediana e a nota mínima foi zero para todas as regiões do Brasil, sem exceção.
Ou seja, pelo menos metade dos alunos tirou nota zero nesta questão.
O Gráfico 3.21 representa a distribuição das notas da questão discursiva 5 do componente de Conhecimento Específico.
Destaca -se, como nas demais questões discursivas, o grande número de estudantes que deixaram esta questão em branco.
O coeficiente de assimetria para todos os alunos é 1,79 e para todas as regiões se mantém positivo, variando entre 1,50 para a região Sul e 2,14 para a região Centro-Oeste. 3.3.2.6 Comentários sobre as respostas à Questão Discursiva 5 O enunciado e a linguagem utilizada não nos pareceram muito claros.
Consideramos que o “ caráter apenas motivador ” do texto de apoio da questão poderia gerar uma expectativa no aluno de que as ideias presentes deveriam ser usadas.
O enunciado não deixava claro se era uma questão para o aluno dar a sua opinião sobre o assunto ou para utilizar fundamentação teórica.
O primeiro ponto percebido nas correções foi que faltou aos estudantes um entendimento do enunciado da questão.
Pouquíssimos concluintes de publicidade demonstraram saber a respeito do movimento da sociedade civil organizada em relação à democratização das comunicações.
Mesmo nas respostas que contemplaram questões como o fim do monopólio na mídia ou a abertura de espaço para iniciativas comunitárias de comunicação, principalmente rádios e TVs, o conteúdo exposto demonstrou ser muito mais parte do senso comum do que do acompanhamento do debate sobre a proposta do marco regulatório.
A maioria das respostas trouxe a ideia de que democratizar a comunicação é garantir que todas as pessoas possam falar o que quiserem, independente de contexto.
Muitos estudantes citaram o período da ditadura militar no Brasil e o direito constitucional de ir e vir para justificar os seus argumentos.
Muitos copiaram literalmente o texto motivador ou repetiram as ideias apresentadas, modificando apenas as palavras.
Alguns mesclaram informações presentes no enunciado das questões anteriores em suas respostas.
Muitos apontaram a Internet como principal instrumento para a democratização da comunicação, principalmente pela facilidade de acesso a partir do barateamento de smartphones, tablets e do próprio sinal 3G.
Mais uma vez, o que mais chamou atenção foi o desrespeito às normas gramaticais e erros grosseiros de ortografia.
Em grande parte dos textos não havia pontuação, principalmente vírgulas.
Muitas respostas pareceram uma colagem de frases desconexas, com falta de clareza nas informações apresentadas.
O último ponto a ressaltar foi o alto índice de respostas em branco, justificado talvez pela extensão da prova, muito citada pelos alunos que fizeram algum tipo de protesto contra o exame.
Alguns começaram a responder à questão e depois sinalizaram que não houve tempo para a conclusão da resposta. 3.3.3 Considerações Finais Ao final do período de correção, alguns pontos que foram observados merecem destaque : • O número de questões em branco foi elevado, podendo ser por protesto ou pelo desconhecimento dos assuntos explorados pela prova. • O desconhecimento dos assuntos também foi muito grande. • Houve graves erros de português.
Porém, o mais preocupante é a incapacidade da grande maioria dos estudantes de se fazer entender e não se mostrar capaz de escrever um texto coerente.
As ideias são muitas vezes confusas, contraditórias e incoerentes.
A capacidade de argumentação é fraca e a dissertação se torna improvável diante deste cenário ( muitas respostas foram entregues em tópicos ).
Além disso, a maior parte das respostas se mostrou opinativa, com uso de verbos em primeira pessoa e o estudante apresentando o que achava do assunto, não necessariamente baseado em teorias ou no aprendizado acadêmico.

% matched lemmas: cada, estar, ordem, torno
\end{textsample}
